\section*{Capítulo \ref{ch:b2:structure-determination} --- Determinación estructural: RMN de \texorpdfstring{$^{13}$C}{13C}, EM e IR combinadas}

\begin{solution}{exo:b2:structure-determination:1}
1,2-Dimetilbenceno: 4 (\ce{CH3}; C1/C2; C3/C6; C4/C5). 1,3-: 5 (\ce{CH3}; C1/C3; C2; C4/C6; C5).
1,4-: 3 (\ce{CH3}; C1/C4; C2/C3/C5/C6).
\end{solution}

\begin{solution}{exo:b2:structure-determination:2}
\ce{CH3CH(OH)CH2CH3}. DEPT-135: C1 (\ce{CH3}) hacia arriba, C2 (CH) hacia arriba, C3 (\ce{CH2}) hacia abajo, C4 (\ce{CH3}) hacia arriba.
DEPT-90: solo C2.
\end{solution}

\begin{solution}{exo:b2:structure-determination:3}
209: carbonilo de cetona; 170: carbonilo de éster (o de ácido); 129: carbono \omterm{def:b2:huckel:aromatic}{aromático}; 64: carbono unido a un
oxígeno; 14: \ce{CH3} alquílico.
\end{solution}

\begin{solution}{exo:b2:structure-determination:4}
$(2 \times 8 + 2 - 8)/2 = 5$: un anillo bencénico (4) y un \ce{C=O} (1). Por ejemplo, el benzoato de metilo,
\ce{C6H5COOCH3}, o el ácido 4-metilbenzoico.
\end{solution}

\begin{solution}{exo:b2:structure-determination:5}
Una insaturación, una cetona (\qty{1715}{cm^{-1}}, línea en 209, cuaternaria); un \ce{CH2} y dos
\ce{CH3}, cuatro carbonos distintos: butan-2-ona, \ce{CH3COCH2CH3}.
\end{solution}

\begin{solution}{exo:b2:structure-determination:6}
Carbono: la pentan-2-ona muestra cinco líneas, y la pentan-3-ona, simétrica, tres. Protones: la pentan-2-ona tiene un
singlete de tres protones (\ce{CH3CO}); la pentan-3-ona, solo un cuadruplete y un triplete. \omterm{def:b2:mass-spec-atomic:spectrum}{Espectro de masas}: la pentan-2-ona
tiene su \omterm{def:b2:mass-spec-atomic:spectrum}{pico base} en 43 (\ce{CH3CO+}) y un ion de McLafferty en 58; la pentan-3-ona tiene su \omterm{def:b2:mass-spec-atomic:spectrum}{pico base} en 57
(\ce{C2H5CO+}).
\end{solution}

\begin{solution}{exo:b2:structure-determination:7}
\ce{C4H8O2}, 88: el par cuadruplete--triplete es un \ce{OCH2CH3} (el \ce{CH2} en 4.12 está sobre oxígeno), y el
singlete, un \ce{CH3CO}. Etanoato de etilo, \ce{CH3COOCH2CH3}.
\end{solution}

\begin{solution}{exo:b2:structure-determination:8}
1,2-Diclorobenceno: tres líneas (C1/C2, C3/C6, C4/C5). 1,4-Diclorobenceno: dos (C1/C4, los cuatro
CH). El DEPT-90 muestra dos líneas de CH para el primero y una para el segundo.
\end{solution}

\begin{solution}{exo:b2:structure-determination:9}
C2 es un estereocentro. Sustituir uno u otro protón de C3 por un deuterio da dos
diastereoisómeros, no enantiómeros: los dos protones son diastereotópicos, en entornos distintos.
Pueden tener desplazamientos distintos y acoplarse entre sí además de con sus vecinos, de modo que el
\ce{CH2} da un multiplete complejo en lugar de un simple quintuplete.
\end{solution}

\begin{solution}{exo:b2:structure-determination:10}
Cinco insaturaciones; un anillo monosustituido (tres líneas de CH y una de C en 137.0, cinco H
\omterm{def:b2:huckel:aromatic}{aromáticos}); una cetona (200.8, C); un grupo etilo (\ce{CH2} 31.8, cuadruplete en 2.98 contiguo al carbonilo, \ce{CH3}
8.3). 1-Fenilpropan-1-ona, \ce{C6H5COCH2CH3}. La conjugación con el anillo deslocaliza los electrones $\pi$ del \ce{C=O}
y debilita el enlace: la banda cae por debajo de los \qty{1715}{cm^{-1}} de una cetona saturada.
% ledger: ir:CO-ketone
\end{solution}

\begin{solution}{exo:b2:structure-determination:11}
Ácido pentanoico: la banda \ce{O-H} muy ancha de 2500 a \qty{3300}{cm^{-1}} y un protón hacia 11--12
ppm. Butanoato de metilo: un singlete de tres protones en el intervalo de los protones de un carbono unido a un oxígeno
(3.3--4.5). Propanoato de etilo: dos cuadrupletes, uno sobre oxígeno y otro contiguo al carbonilo. Etanoato de
propilo: un singlete de tres protones contiguo al carbonilo (2.0--2.4) y un triplete sobre oxígeno.
% ledger: ir:OH-acid, nmrr:acid, nmrr:alcohol, nmrr:ketone
\end{solution}

\begin{solution}{exo:b2:structure-determination:12}
Ambas líneas desaparecen en \omterm{def:b2:structure-determination:dept}{DEPT}, de modo que ambas son cuaternarias, y el \omterm{def:b2:structure-determination:dept}{DEPT} por sí solo no puede distinguirlas. Los desplazamientos
sí: los carbonilos de éster se sitúan hacia 167--171 y los carbonos \omterm{def:b2:huckel:aromatic}{aromáticos} entre 128 y 137 en el gráfico de este
capítulo. Un carbono del anillo en 166.5 y un carbonilo en 130.6 quedarían ambos muy fuera de sus clases; la
asignación correcta es la inversa.
\end{solution}

\begin{solution}{pb:b2:structure-determination:1}
\textbf{1.} $3.1/31.7 \approx 9.8~\%$; $9.8/1.11 \approx 9$ carbonos.
\textbf{2.} Masa par: ningún nitrógeno (o dos).
\textbf{3.} $150 - 108 = 42$: \ce{H10O2} encaja (10 + 32): \ce{C9H10O2}.
\textbf{4.} Diez carbonos darían $M+1 \approx 11~\%$ de $M$, no 9.8~\%; y el IR muestra un
carbonilo, que necesita un segundo oxígeno además de la banda \ce{C-O}.
\textbf{5.} $(2 \times 9 + 2 - 10)/2 = 5$.
\textbf{6.} $108 + 10 \times 1.007825 + 2 \times 15.994915 \approx 150.0681$.
\textbf{7.} Una tensión \ce{C=O} en la posición de un éster saturado (hacia \qty{1735}{cm^{-1}}).
\textbf{8.} La tensión \ce{C-O} del éster.
\textbf{9.} Enlaces \ce{C-H} \omterm{def:b2:huckel:aromatic}{aromáticos} (o alquénicos).
\textbf{10.} Un grupo \ce{O-H}: ni alcohol ni ácido carboxílico.
\textbf{11.} No: un carbonilo conjugado con un anillo absorbe a menor número de onda; aquí se sitúa donde lo hacen
los ésteres saturados, de modo que un grupo no conjugante lo separa del anillo.
\textbf{12.} 170.70: \ce{C=O} de éster; 136.14: carbono del anillo que lleva el sustituyente; 128.56 y 128.24:
CH \omterm{def:b2:huckel:aromatic}{aromáticos}; 66.24: \ce{CH2} unido a un oxígeno; 20.82: \ce{CH3}.
\textbf{13.} 66.24, el \ce{CH2}.
\textbf{14.} 170.70 y 136.14, los dos carbonos sin hidrógeno.
\textbf{15.} 128.56 y 128.24.
\textbf{16.} Cuatro: el carbono que lleva el sustituyente, los dos orto, los dos meta y el CH para. Se esperan tres
tipos de CH; dos de ellos tienen desplazamientos demasiado próximos para resolverse a este campo, y comparten una línea.
\textbf{17.} No necesariamente: las alturas no son recuentos (\cref{prop:b2:structure-determination:intensities}),
aunque aquí la línea bien puede contener dos tipos de CH solapados.
\textbf{18.} 7.33: los cinco protones \omterm{def:b2:huckel:aromatic}{aromáticos}; 5.09: el \ce{OCH2}; 2.06: el \ce{CH3CO}.
\textbf{19.} Sus átomos vecinos (el carbono del anillo y el oxígeno para el \ce{CH2}, el carbono carbonílico
para el \ce{CH3}) no llevan hidrógeno.
\textbf{20.} Está unido al oxígeno electronegativo del éster y al anillo: ambos lo desapantallan.
\textbf{21.} \ce{C6H5CH2OCOCH3}, etanoato de bencilo.
\textbf{22.} En el isómero el \ce{CH2} no está unido a un oxígeno, de modo que no podría aparecer en 5.09, y el
metilo, sobre oxígeno, aparecería en el intervalo 3.3--4.5, no en 2.06 junto a un carbonilo.
% ledger: nmrr:alcohol, nmrr:ketone
\textbf{23.} 108: pérdida de 42 (cetena, \ce{CH2=C=O}) por transposición, el catión radical de
\ce{C7H8O}; 91: \ce{C7H7+}, el catión bencilo; 43: \ce{CH3CO+}.
\textbf{24.} Siete tipos de carbono (\ce{C=O}, el carbono sustituido del anillo, CH orto, meta y para,
\ce{CH2}, \ce{CH3}): \textbf{siete líneas, una de ellas (la del \ce{CH2}) hacia abajo} en DEPT-135.
\end{solution}
